%% RevTeX 4-2 paper --- compiled for APS / Physical Review A
\documentclass[aps,pra,reprint,superscriptaddress,amsmath,amssymb]{revtex4-2}

\usepackage{graphicx}
\usepackage{xcolor}
\usepackage{amsthm}
\usepackage[expansion=false]{microtype}
\usepackage{enumitem}
\usepackage{hyperref}
\usepackage{orcidlink}

\hypersetup{
    colorlinks=true,
    linkcolor=blue,
    filecolor=magenta,
    urlcolor=blue,
    citecolor=blue,
    pdftitle={Against a Universal Trading Strategy: No-Arbitrage, No-Free-Lunch, and Adversarial Counterpaths},
    pdfauthor={Karl Svozil},
    pdfsubject={Mathematical limits on universally profitable trading strategies},
    pdfkeywords={no-arbitrage, no-free-lunch, martingales, computability, algorithmic randomness, universal portfolios}
}

%% Mathematical Environments
\theoremstyle{plain}
\newtheorem{theorem}{Theorem}
\newtheorem{lemma}[theorem]{Lemma}
\newtheorem{corollary}[theorem]{Corollary}
\newtheorem{proposition}[theorem]{Proposition}
\theoremstyle{definition}
\newtheorem{definition}{Definition}
\newtheorem{criterion}{Criterion}
\theoremstyle{remark}
\newtheorem{remark}{Remark}

\begin{document}

\title{Against a Universal Trading Strategy: No-Arbitrage, No-Free-Lunch, and Adversarial Counterpaths}

\author{Karl Svozil\,\orcidlink{0000-0001-6554-2802}}
\affiliation{Institute for Theoretical Physics, TU Wien, Wiedner Hauptstrasse 8-10/136, A-1040 Vienna, Austria}
\email{karl.svozil@tuwien.ac.at}
\homepage{http://tph.tuwien.ac.at/~svozil}

\date{July 27, 2026}

\begin{abstract}
We investigate universally winning trading strategies, defined as admissible,
self-financing strategies whose discounted gain from zero initial capital is
strictly positive on every admissible market trajectory. Such a strategy is a
classical arbitrage, and its pathwise requirement is stronger than the usual
almost-sure condition. It is therefore excluded already by no-arbitrage and,
a fortiori, in locally bounded semimartingale markets satisfying No Free Lunch
with Vanishing Risk and admitting an equivalent local martingale measure.
Separate modeling regimes yield complementary, but logically distinct,
limitations. In a uniform binary prediction model, every causal predictor has
mean directional accuracy \(1/2\); optional stopping nevertheless permits an
arbitrarily high trade win rate while preserving zero expected gain. In a
worst-case computational model, every total deterministic trading algorithm
has a computable adaptive counterpath on which its gain is non-positive.
Ville's inequality and computable randomness connect the probabilistic and
computational regimes without requiring an adaptive adversary. We make the
thermodynamic analogy precise only at the level of positive density
martingales and fluctuation identities, rather than identifying a martingale
measure with detailed balance or invoking financially relevant Landauer
costs. Finally, an explicit put--call-parity analysis of the Wheel Options
Strategy shows that its apparent income is compensation for put-like downside
and capped upside. Positive results such as Cover's universal portfolio are
shown to concern relative regret, not guaranteed absolute profit.
\end{abstract}

\maketitle

%%------------------------------------------------------------------------------
\section{Introduction}
\label{sec:intro}
%%------------------------------------------------------------------------------

The question of whether a trading strategy can always make money, regardless
of what the market does, is older than modern mathematical finance. The
mathematical study of unpredictable prices begins with Bachelier's 1900
thesis~\cite{Bachelier1900}, continues through Samuelson's martingale
formulation~\cite{Samuelson1965}, and informs Fama's Efficient Market
Hypothesis (EMH)~\cite{Fama1970}. The EMH is an economic claim about the
incorporation of information into prices; it is not, by itself, a theorem that
no individual strategy can win unconditionally.

Here a strategy is called \emph{universal} only if, after discounting by a
specified numeraire, it produces a strictly positive terminal gain from zero
initial capital on every admissible market trajectory. The zero-cost,
self-financing, and discounted-gain qualifications are essential: a positive
initial investment in a bank account has positive terminal wealth, but it is
not an arbitrage and does not outperform its financing benchmark.

This paper articulates the resulting impossibility through three distinct
theoretical paradigms, with their assumptions kept separate:

\begin{enumerate}
  \item \textbf{Financial (Measure-Theoretic, Sec.~\ref{sec:noarb}).}
  Pathwise universality is a special case of classical arbitrage. It is
  excluded by no-arbitrage and, under standard semimartingale hypotheses, by
  the existence of an equivalent local martingale measure. Optional stopping
  also explains why a high fraction of profitable trades need not imply
  positive expected gain.

  \item \textbf{Combinatorial (Uniform Prediction, Sec.~\ref{sec:nfl}).}
  The Wolpert--Macready No-Free-Lunch (NFL) theorem concerns uniform averaging
  over finite objective-function spaces~\cite{Wolpert:1997:NFL:2221336.2221408}.
  For trading-direction claims we prove the relevant statement directly:
  averaged uniformly over all binary price-direction sequences, every causal
  predictor has accuracy \(1/2\).

  \item \textbf{Computational (Adaptive Counterpaths,
  Sec.~\ref{sec:computability}).} Against every total deterministic computable
  strategy, a computable environment can simulate its announced position and
  choose the next return with the opposite sign. A separate
  martingale-randomness result describes what happens under a passive
  computable fair measure.
\end{enumerate}

We frame these findings with a \emph{time-reversal stress test}
(Sec.~\ref{sec:criterion}) and a carefully limited Maxwell's Demon analogy
(Sec.~\ref{sec:demon}). The social-choice comparison in
Sec.~\ref{sec:arrow} is explicitly analogical rather than a reduction of one
theorem to another. The Wheel Options Strategy is analyzed in
Sec.~\ref{sec:wheel}, Cover's different notion of universality is separated
from absolute profit in Sec.~\ref{sec:cover}, and the conditional scope of
strategies working \emph{for all practical purposes} (FAPP) is discussed in
Sec.~\ref{sec:fapp}.

%%------------------------------------------------------------------------------
\section{The Time-Reversal Stress Test}
\label{sec:criterion}
%%------------------------------------------------------------------------------

Let \(\mathcal{M}\) be a space of strictly positive paths
\(P\colon[0,T]\to\mathbb{R}_{>0}^n\). The time reversal of \(P\) is
\begin{equation}
  \widetilde{P}(t)=P(T-t), \qquad t\in[0,T].
  \label{eq:timereversal}
\end{equation}
For a causal strategy \(\sigma\), let \(\Pi_T[\sigma;P]\) denote its terminal
discounted gain from zero initial capital when evaluated on \(P\).

\begin{definition}[Pathwise Universal Strategy]
\label{def:universal}
  A causal strategy \(\sigma\) is \emph{pathwise universal} on
  \(\mathcal{M}\) if it is self-financing and admissible and satisfies
  \(\Pi_T[\sigma;P]>0\) for every \(P\in\mathcal{M}\).
\end{definition}

\begin{criterion}[Reversal Falsification Test]
  \label{crit:tr}
  Suppose \(\mathcal{M}\) is closed under time reversal. If there exists
  \(P\in\mathcal{M}\) such that
  \begin{equation}
    \min\{\Pi_T[\sigma;P],\Pi_T[\sigma;\widetilde P]\}\leq 0,
    \label{eq:reversaltest}
  \end{equation}
  then \(\sigma\) is not pathwise universal on \(\mathcal{M}\).
\end{criterion}

\begin{proof}
Both \(P\) and \(\widetilde P\) belong to \(\mathcal{M}\). A non-positive gain
on either member of the pair contradicts Definition~\ref{def:universal}.
\end{proof}

Criterion~\ref{crit:tr} is deliberately a falsification device, not an
independent impossibility theorem. Reversing a profitable path does not in
general negate the profit of a causal strategy, because the strategy can take
different positions on the reversed history. The value of the test is that it
generates a paired stress scenario whenever the modeled path class is
reversal-closed.

\begin{remark}[Filtrations and Time Reversal]
In stochastic finance, the set of admissible trajectories is not automatically
closed under time reversal. The process \(P(T-t)\) is generally not adapted to
the original forward filtration, and its semimartingale property under a
reversed filtration requires additional hypotheses. The stress test is
therefore pathwise; the rigorous probabilistic constraint is given next.
\end{remark}

%%------------------------------------------------------------------------------
\section{The No-Arbitrage Foundation}
\label{sec:noarb}
%%------------------------------------------------------------------------------

\subsection{Discounted Gains and Admissibility}

Let the market be defined on a filtered probability space
\((\Omega,\mathcal{F},\{\mathcal{F}_t\}_{t\in[0,T]},\mathbb{P})\).
Let \(B_t>0\) be a numeraire and \(S_t\) the vector of traded asset prices.
Write \(\overline S_t=S_t/B_t\) for discounted prices. A predictable,
\(\overline S\)-integrable strategy \(H\) has discounted self-financing gain
\begin{equation}
  G_t=(H\mathbin{\cdot}\overline S)_t,\qquad G_0=0.
  \label{eq:gain}
\end{equation}
It is \emph{admissible} if there is a finite constant \(a\) such that
\begin{equation}
  G_t\geq -a\quad\text{for all }t\in[0,T],\quad\mathbb{P}\text{-a.s.}
  \label{eq:admissible}
\end{equation}
This uniform lower bound rules out doubling systems financed by unbounded
credit.

\begin{theorem}[Universality Implies Classical Arbitrage]
  \label{thm:arbitrage}
  If an admissible self-financing strategy \(H\) with zero initial capital
  satisfies \(G_T(\omega)>0\) for every \(\omega\in\Omega\), then \(H\) is a
  classical arbitrage.
\end{theorem}

\begin{proof}
The pathwise assumption implies \(G_T\geq0\) almost surely and
\(\mathbb{P}(G_T>0)=1\). Together with zero initial cost, self-financing, and
admissibility, these are stronger than the usual classical-arbitrage
conditions \(\mathbb{P}(G_T\geq0)=1\) and
\(\mathbb{P}(G_T>0)>0\).
\end{proof}

\begin{remark}[Terminology and Strength]
The pathwise condition is sometimes called a sure or strong arbitrage, but it
is not synonymous with \emph{arbitrage of the first kind}. The latter is the
condition dual to No Unbounded Profit with Bounded Risk (NUPBR/NA1) and is a
distinct, weaker viability concept~\cite{Kardaras2012}. Classical
no-arbitrage alone already excludes the universal strategy in
Theorem~\ref{thm:arbitrage}; NFLVR is invoked below because it supplies the
standard martingale-measure formulation.
\end{remark}

\subsection{The Martingale Constraint}

For a locally bounded semimartingale price process, the
Delbaen--Schachermayer form of the First Fundamental Theorem states that No
Free Lunch with Vanishing Risk (NFLVR) is equivalent to the existence of an
equivalent local martingale measure (ELMM)
\(\mathbb{Q}\approx\mathbb{P}\) under which \(\overline S\) is a local
martingale~\cite{Harrison1981,Delbaen1994}.

\begin{corollary}[No Universal Zero-Cost Strategy]
  \label{cor:noarb}
  If an ELMM \(\mathbb{Q}\) exists, no admissible self-financing strategy with
  zero initial capital can have \(G_T>0\) on every path.
\end{corollary}

\begin{proof}
Under \(\mathbb{Q}\), the admissible stochastic integral
\(G=H\mathbin{\cdot}\overline S\) is a local martingale bounded from below and
therefore a supermartingale. Hence
\(\mathbb{E}_{\mathbb{Q}}[G_T]\leq G_0=0\). If \(G_T>0\) on every path, then
equivalence gives \(G_T>0\), \(\mathbb{Q}\)-a.s., and consequently
\(\mathbb{E}_{\mathbb{Q}}[G_T]>0\), a contradiction.
\end{proof}

The ELMM conclusion concerns discounted gains under the pricing measure
\(\mathbb{Q}\). It does not say that every risky strategy has non-positive
expected return under the physical measure \(\mathbb{P}\); positive physical
expected returns may compensate risk.

\subsection{Win Rate Is Not Expected Value}

The probability of closing a trade at a profit can greatly exceed \(1/2\)
even in a fair market. Optional stopping gives the precise distinction.

\begin{proposition}[Asymmetric Barriers in a Fair Random Walk]
  \label{prop:optional}
  Let \(X_n\) be a simple symmetric random walk with \(X_0=0\), and for positive
  integers \(a,b\) let
  \[
    \tau=\inf\{n\geq0:X_n=a\text{ or }X_n=-b\}.
  \]
  Then
  \begin{equation}
    \mathbb{P}(X_\tau=a)=\frac{b}{a+b},
    \qquad
    \mathbb{E}[X_\tau]=0.
    \label{eq:gambler}
  \end{equation}
\end{proposition}

\begin{proof}
The stopped process \(X_{n\wedge\tau}\) is a bounded martingale, so Doob's
optional-stopping theorem applies~\cite{Doob1953}. If
\(p=\mathbb{P}(X_\tau=a)\), then
\(0=\mathbb{E}[X_\tau]=ap-b(1-p)\), which gives
\(p=b/(a+b)\).
\end{proof}

With a take-profit of \(a=1\) and a stop-loss of \(b=9\), the win probability
is \(0.9\), but the expected payoff is zero: nine frequent unit gains are
offset by one nine-unit loss in expectation. Thus prediction accuracy, trade
win rate, and expected discounted profit are three different quantities.

%%------------------------------------------------------------------------------
\section{The Maxwell's Demon Analogy: A Density-Martingale Boundary}
\label{sec:demon}
%%------------------------------------------------------------------------------

\subsection{Formal Analogy, Not Equivalence}

Maxwell's Demon observes microscopic degrees of freedom and uses information
to extract work from fluctuations. A trading algorithm observes a signal and
uses it to choose a position. The comparison is useful only after the
equilibrium constraints on both systems are stated explicitly
(Table~\ref{tab:demon}).

\begin{table}[htbp]
\caption{\label{tab:demon}A limited analogy between information engines and trading strategies.}
\centering
\begin{ruledtabular}
\begin{tabular}{p{0.42\columnwidth}p{0.48\columnwidth}}
\textbf{Information engine} & \textbf{Trading strategy} \\
\colrule
Observes microscopic data & Observes prices or signals \\
Uses feedback to extract work & Uses positions to seek gain \\
Requires cyclic accounting & Requires self-financing accounting \\
Bounded by entropy/information laws & Bounded by no-arbitrage under an ELMM \\
\end{tabular}
\end{ruledtabular}
\end{table}

There is a precise martingale identity behind the analogy. If
\(\mathbb{Q}\approx\mathbb{P}\), define the density process
\begin{equation}
  Z_t=\mathbb{E}_{\mathbb{P}}\!\left[
    \left.\frac{d\mathbb{Q}}{d\mathbb{P}}\right|\mathcal{F}_t
  \right].
  \label{eq:density}
\end{equation}
Then \(Z\) is a positive \(\mathbb{P}\)-martingale with
\(\mathbb{E}_{\mathbb{P}}[Z_T]=1\). Writing
\(\Sigma_T=-\log Z_T\) gives the fluctuation-form identity
\begin{equation}
  \mathbb{E}_{\mathbb{P}}\!\left[e^{-\Sigma_T}\right]=1.
  \label{eq:fluctuation}
\end{equation}
Equation~\eqref{eq:fluctuation} has the same normalization structure as an
integral fluctuation relation such as Jarzynski's
equality~\cite{Jarzynski1997}. Here, however, \(\Sigma_T\) is a
log-likelihood ratio between two financial measures, not physical entropy
production.

An ELMM is also not the financial equivalent of thermodynamic detailed
balance. The martingale property imposes zero conditional drift on discounted
traded prices under \(\mathbb{Q}\); detailed balance additionally imposes a
time-reversal symmetry on a stationary transition law. The latter is strictly
stronger.

\subsection{Information Has a Financial Value, but Landauer Cost Does Not}

Landauer's bound for erasing one bit at room temperature is approximately
\(k_BT\ln2\approx2.9\times10^{-21}\,\mathrm{J}\)
per bit~\cite{landauer:61,bennett-82}. Converting this into a ``cost per
trade'' requires both the number of logically erased bits and an energy-price
model. Even then, the bound is quantitatively negligible compared with
spreads, impact, latency, and model risk; it is not the mechanism enforcing
financial no-arbitrage.

The economically relevant information bound is instead expressed in growth
units. In the canonical Kelly horse-race model with fair odds, the increase in
optimal expected log wealth produced by side information \(Y\) about outcome
\(X\) is the mutual information \(I(X;Y)\)~\cite{Kelly1956}. This result is
model-specific, but it makes the demon analogy quantitative without confusing
thermodynamic energy with money: exploitable information can increase optimal
growth only because the joint law of signal and outcome is non-uniform.

%%------------------------------------------------------------------------------
\section{The Combinatorial Perspective}
\label{sec:nfl}
%%------------------------------------------------------------------------------

The optimization NFL theorem and the elementary prediction result needed for
market-direction claims should be distinguished.

\begin{theorem}[Optimization NFL, Wolpert--Macready]
  \label{thm:nfl}
  For finite search and cost spaces, when performance is averaged uniformly
  over all objective functions, any two non-revisiting black-box search
  algorithms have identical averaged performance.
\end{theorem}

The theorem concerns oracle-based optimization over function spaces. It does
not, by itself, imply a statement about self-financing gains on stochastic
price paths. For binary directional prediction the relevant result is simpler
and can be derived directly.

\begin{proposition}[Uniform Binary Prediction]
  \label{prop:binary}
  Let \(X_1,\ldots,X_T\) be uniformly distributed on \(\{0,1\}^T\). A causal
  deterministic predictor chooses
  \(\widehat X_t=f_t(X_1,\ldots,X_{t-1})\). Its mean accuracy
  \[
    A_T=\frac1T\sum_{t=1}^T
      \mathbf{1}\{\widehat X_t=X_t\}
  \]
  satisfies \(\mathbb{E}[A_T]=1/2\). The same holds for a randomized predictor
  whose private randomization is independent of the next bit.
\end{proposition}

\begin{proof}
Conditional on every history \(X_1,\ldots,X_{t-1}\), the next bit \(X_t\) is
an independent fair bit. Therefore
\(\mathbb{P}(\widehat X_t=X_t\mid X_1,\ldots,X_{t-1})=1/2\).
Linearity of expectation gives the result. Conditioning additionally on the
predictor's private randomization proves the randomized case.
\end{proof}

Consequently, no predictor can achieve expected directional accuracy greater
than \(1/2\) under the uniform distribution over all binary sequences; in
particular, none can exceed \(1/2\) on every sequence. This is an NFL symmetry
statement, not a model of actual markets. Financial sign sequences are not
uniform in general, and a directional forecast does not determine trading
profit because payoff sizes, financing, and stopping rules matter
(Proposition~\ref{prop:optional}).

The defensible practical conclusion is conditional: every successful
predictor or strategy exploits non-uniform structure in the physical measure
\(\mathbb{P}\). Its edge is evidence about a restricted environment, not
universal algorithmic superiority.

%%------------------------------------------------------------------------------
\section{The Computational Perspective: Adaptive Counterpaths}
\label{sec:computability}
%%------------------------------------------------------------------------------

We now discard the assumption of an exogenous price process and let the market
react after observing the trader's current position. This is a worst-case,
perfect-information model. Let a trading program \(M_\sigma\) be a
\emph{total} computable map from every finite computable price history \(H_t\)
to a rational position \(w_t\). Totality models the practical requirement that
a deployed algorithm return an action before a deadline.

\begin{theorem}[Computable Adaptive Counterpath]
  \label{thm:counterpath}
  For every total deterministic computable strategy \(M_\sigma\) and every
  finite horizon \(T\), there exists a positive computable price path
  \(P_{\mathrm{adv}}\) on which its self-financing gain
  \[
    G_T=\sum_{t=0}^{T-1}w_t(P_{t+1}-P_t)
  \]
  is non-positive. It is strictly negative if the strategy takes a nonzero
  position at least once.
\end{theorem}

\begin{proof}
Choose computable \(P_0>0\) and rational \(0<\epsilon<1\). At time \(t\),
compute \(w_t=M_\sigma(H_t)\) and set
\begin{equation}
P_{t+1}=
\begin{cases}
P_t(1-\epsilon),&w_t>0,\\
P_t(1+\epsilon),&w_t<0,\\
P_t,&w_t=0.
\end{cases}
\label{eq:adversary}
\end{equation}
The resulting prices remain positive and computable. In each period,
\(w_t(P_{t+1}-P_t)\leq0\), with strict inequality when \(w_t\neq0\).
Summing proves the claim.
\end{proof}

This is a simulate-and-oppose construction, not a reduction from the Halting
Problem. Its quantifiers are
\(\forall M_\sigma\,\exists P_{\mathrm{adv}}\): the counterpath may depend on
the strategy. There need not be one finite path on which every strategy loses;
for example, a permanently long and a permanently short position cannot both
lose on the same one-step price move.

\begin{corollary}[Maximin Value]
  \label{cor:maximin}
  In the finite-horizon game of Theorem~\ref{thm:counterpath}, if the neutral
  strategy is allowed, then
  \[
    \sup_{M_\sigma}\inf_{P_{\mathrm{adv}}}G_T(M_\sigma,P_{\mathrm{adv}})=0.
  \]
\end{corollary}

\begin{proof}
The neutral strategy guarantees zero, while
Theorem~\ref{thm:counterpath} gives every strategy a counterpath with gain at
most zero.
\end{proof}

\begin{proposition}[Private Randomization Does Not Create a Guarantee]
  \label{prop:randomized}
  Suppose a randomized trader's current coin toss is hidden from the
  environment. There still exists a passive computable distribution of next
  returns under which every integrable strategy has conditional expected gain
  zero.
\end{proposition}

\begin{proof}
Let the next price increment be \(+\epsilon P_t\) or \(-\epsilon P_t\) with
equal probability, independently of the trader's private draw. Conditional on
the current information and position,
\(\mathbb{E}[w_t(P_{t+1}-P_t)\mid\mathcal{F}_t]=0\).
Summing gives expected gain zero.
\end{proof}

Randomization therefore prevents a code-observing opponent from mechanically
opposing the unseen current draw, but it does not yield strictly positive
expected gain in every environment.

\subsection{The Algorithmic-Randomness Bridge}

There is also a non-adaptive connection between computability and martingale
finance. In the fair binary market, a nonnegative capital process
\(K\colon\{0,1\}^{<\mathbb{N}}\to\mathbb{R}_{\geq0}\) is a martingale when
\begin{equation}
  K(s)=\frac{K(s0)+K(s1)}2.
  \label{eq:computablemartingale}
\end{equation}

\begin{theorem}[Ville's Inequality]
  \label{thm:ville}
  If \(K_0=1\) is a nonnegative martingale under the fair-coin measure, then
  for every \(\lambda>0\),
  \begin{equation}
    \mathbb{P}\!\left(\sup_{n\geq0}K_n\geq\lambda\right)
    \leq\frac1\lambda.
    \label{eq:ville}
  \end{equation}
\end{theorem}

Ville's inequality follows by stopping the nonnegative capital process when it
first crosses \(\lambda\) and applying optional stopping, followed by a
limit~\cite{Ville1939}. In particular, the set on which a fixed nonnegative
martingale becomes unbounded has probability zero.

A binary sequence is \emph{computably random} precisely when no computable
nonnegative martingale succeeds on it by becoming unbounded~\cite{schnorr1,DH}.
Since there are only countably many computable
martingales, almost every fair-coin sequence is computably random. Thus a
single typical passive path prevents unbounded enrichment by every computable
fair-game betting system. This statement differs from
Theorem~\ref{thm:counterpath}: it is asymptotic, measure-one, and concerns
unbounded success, whereas the adaptive counterpath gives a finite
non-positive gain tailored to one strategy. Game-theoretic probability
develops this pathwise betting viewpoint systematically~\cite{ShaferVovk2001}.

\subsection{What Rice's Theorem Does and Does Not Say}

Rice's theorem states that every nontrivial extensional property of arbitrary
partial computable functions is undecidable~\cite{Rice-1953}. It therefore
limits fully general program certification; even deciding whether arbitrary
code is total is impossible in general~\cite{turing-36}. It does \emph{not}
imply that the loss-producing inputs of a fixed total finite-horizon strategy
are undecidable. On a specified finite computable path, such a program can be
simulated and its gain calculated. No probability claim about market failure
sets follows from Rice's theorem alone.

%%------------------------------------------------------------------------------
\section{Analogy to Social Choice: Scope and Limits}
\label{sec:arrow}
%%------------------------------------------------------------------------------

Arrow's Impossibility Theorem concerns social welfare functions on at least
three alternatives. Under unrestricted preference profiles, collective
rationality, Pareto unanimity, and independence of irrelevant alternatives,
the aggregation rule must be dictatorial~\cite{Arrow1963}. Its assumptions and
objects differ from those of a trading model.

There is nevertheless a useful structural comparison:
\begin{description}
  \item[Universal domain.] Arrow quantifies over all admissible preference
  profiles; pathwise trading universality quantifies over all admissible price
  paths.
  \item[Jointly incompatible demands.] Arrow's fairness and consistency
  requirements cannot all be met by a non-dictatorial rule. In trading,
  zero-cost strict profit on every path is incompatible with no-arbitrage.
  \item[No theorem transfer.] Neither result proves the other. The parallel is
  a shared warning that unrestricted domains make apparently natural
  universal requirements inconsistent.
\end{description}

Yanofsky shows how many self-referential arguments can be organized using
fixed points and fixed-point-free transformations~\cite{Yanofsky-BSL:9051621}.
The response in Eq.~\eqref{eq:adversary} anti-aligns returns with every
nonzero position, but it admits the neutral fixed point \(w_t=0\). That fixed
point yields zero gain rather than strict profit. The relevant obstruction is
therefore the absence of a \emph{strictly profitable} fixed response, not the
absence of every fixed point.

%%------------------------------------------------------------------------------
\section{Case Study: The Wheel Options Strategy}
\label{sec:wheel}
%%------------------------------------------------------------------------------

The Wheel sells a cash-secured put and, after assignment, sells a covered call
against the acquired stock. Its economic exposure follows directly from
put--call parity. Ignoring dividends and financing for the terminal-payoff
identity,
\begin{equation}
  S_T-(S_T-K)^+
  =K-(K-S_T)^+.
  \label{eq:putcall}
\end{equation}
The left-hand side is a covered-call payoff; the right-hand side is cash
paying \(K\) plus a short put. After the corresponding time-zero costs are
matched, a covered call and a cash-secured put are equivalent European
payoffs. Rolling the Wheel therefore repeatedly maintains a put-like,
short-downside-convexity exposure rather than creating directionless income.

\subsection{Failure I: Catastrophic Drop}

If the underlying crashes after put assignment, the short call liability
decreases, which benefits the call writer, but the loss on the stock dominates
and the combined covered-call position can suffer a large net loss.
Equation~\eqref{eq:putcall} makes this explicit: below \(K\), the payoff falls
one-for-one with \(S_T\). A reversed rally/crash pair can be used as the stress
test in Criterion~\ref{crit:tr}, but time reversal alone is not the proof; the
payoff identity is.

\subsection{Failure II: Persistent Decline}

Repeated option premiums offset part of a persistent decline but do not remove
the renewed short-put exposure. The loss is path-dependent rather than
generally linear: strikes, maturities, volatility, assignment dates, and
position sizing determine the sequence of gains.

\subsection{Failure III: Violent Breakout and Benchmark Underperformance}

In a large rally, shares may be called away at the call strike. The Wheel can
still realize a positive absolute profit, so opportunity cost alone is not a
counterexample to strict positive P\&L. It is instead a failure relative to a
buy-and-hold benchmark: Eq.~\eqref{eq:putcall} caps the terminal payoff at
\(K\). No-arbitrage constrains the price of this payoff under
\(\mathbb{Q}\); it does not say that its physical expected return must be zero.

\subsection{Failure IV: Absorbing Ruin Under Capital Constraints}

A fully cash-secured, unlevered cycle has bounded loss, but an underlying that
falls to zero can consume almost all committed capital net of premium.
Leverage makes the absorbing nature of ruin explicit. If wealth evolves as
\(W_{t+1}=W_t(1+fR_{t+1})\) and an admissible return satisfies
\(1+fR_{t+1}=0\), then \(W_{t+1}=0\); a self-financing strategy cannot recover
without an external capital injection. This is a wealth-dynamics fact, not a
consequence of undecidability.

\begin{table*}[ht]
\caption{\label{tab:failuremodes}Wheel failure modes and the mathematical mechanism illustrating each one.}
\begin{ruledtabular}
\begin{tabular*}{\textwidth}{@{\extracolsep{\fill}}llll}
\textbf{Failure} & \textbf{Trajectory Class} & \textbf{Relevant Result} & \textbf{Mechanism} \\
\colrule
I: Crash & Large downside move & Eq.~\eqref{eq:putcall} & Put-like downside \\
II: Bleed & Persistent negative drift & Repeated payoff exposure & Premiums do not insure the stock \\
III: Breakout & Large upside move & Eq.~\eqref{eq:putcall} & Capped benchmark return \\
IV: Ruin & Capital exhaustion & Wealth recursion & Zero wealth is absorbing \\
\end{tabular*}
\end{ruledtabular}
\end{table*}

Table~\ref{tab:failuremodes} is a payoff-based classification, not a claim
that each trajectory has positive probability under every continuous price
model. Corollary~\ref{cor:noarb} excludes strict positive discounted gain on
all paths, but probabilities of particular failure classes require a specified
physical measure.

%%------------------------------------------------------------------------------
\section{Universal Portfolios Are Relative}
\label{sec:cover}
%%------------------------------------------------------------------------------

The word ``universal'' also appears in a celebrated positive theorem.
For price-relative vectors \(x_t\in\mathbb{R}_{>0}^m\), a constant-rebalanced
portfolio \(b\) has wealth
\[
  S_T(b)=\prod_{t=1}^T b^\mathsf{T}x_t.
\]
Cover's universal portfolio is nonanticipating and, under the conditions of
his theorem, has the same asymptotic exponential growth rate as the best
constant-rebalanced portfolio chosen in hindsight on every individual bounded
sequence~\cite{Cover1991}. In regret notation,
\begin{equation}
  \frac1T\left[
  \log\max_{b}S_T(b)-\log S_T^{\mathrm{UP}}
  \right]\longrightarrow0.
  \label{eq:cover}
\end{equation}

Equation~\eqref{eq:cover} is a relative guarantee, not an absolute-profit
guarantee. If every constant-rebalanced portfolio loses money on a path, the
universal portfolio may lose as well while retaining vanishing per-period
regret. Cover's result therefore does not contradict
Corollary~\ref{cor:noarb}; it illustrates the strongest useful meaning of
``universal'' once the benchmark class is specified.

%%------------------------------------------------------------------------------
\section{FAPP Strategies and Their Honest Scope}
\label{sec:fapp}
%%------------------------------------------------------------------------------

While pathwise universal profit is precluded, strategies may be profitable
\emph{for all practical purposes} (FAPP, following Bell~\cite{bell-a})
relative to a physical measure \(\mathbb{P}\), a benchmark, and a finite
horizon. Market making may earn spreads when adverse selection and inventory
risk are controlled, and Kelly investing maximizes expected logarithmic
growth when the relevant distribution is sufficiently well specified. Neither
is an unconditional profit guarantee.

Economic equilibrium also leaves room for conditional edges. Grossman and
Stiglitz show that perfectly informationally efficient markets are
inconsistent with costly information acquisition: if prices revealed all
information without reward, agents would not pay to produce that
information~\cite{GrossmanStiglitz1980}. Conversely, the Milgrom--Stokey
no-trade theorem shows that, under common priors, common knowledge of
rationality, risk aversion, and an initially Pareto-optimal allocation, private
information alone cannot induce mutually agreeable non-null
trade~\cite{MilgromStokey1982}. These are equilibrium boundaries, not
substitutes for the FTAP.

Lo's Adaptive Markets Hypothesis~\cite{Lo2004} supplies a less idealized
interpretation. As capital enters a successful strategy, competition, price
impact, and strategic response may erode its edge. The market need not become
the perfect adversary of Theorem~\ref{thm:counterpath}; adaptation merely makes
the physical distribution endogenous to widespread strategy use.

Automated execution can amplify feedback when many systems share triggers,
funding constraints, or risk models. The 1987 portfolio-insurance
cascade~\cite{Brady1988} and the 2010 Flash Crash~\cite{Kirilenko2017} illustrate
such interactions. Their relevance comes from market impact and coupled
control rules, not from Rice's theorem.

%%------------------------------------------------------------------------------
\section{Conclusion}
\label{sec:conclusion}
%%------------------------------------------------------------------------------

The impossibility of a universally winning zero-cost trading strategy is a
direct consequence of no-arbitrage. Once gains are measured relative to a
numeraire and admissibility and self-financing are stated explicitly,
pathwise strict profit implies classical arbitrage. In a locally bounded
semimartingale market satisfying NFLVR, an equivalent local martingale measure
provides the standard expectation-based contradiction.

The complementary paradigms answer different questions. Uniform binary
averaging fixes directional prediction accuracy at \(1/2\), while optional
stopping shows that trade win rate can nevertheless be arbitrarily high
without positive expected value. An adaptive computable environment can
construct a finite counterpath against each total deterministic strategy.
Under a passive fair measure, Ville's inequality and computable randomness
show instead that typical paths prevent unbounded success by every computable
martingale. None of these statements should be substituted for another.

The thermodynamic comparison is rigorous at the level of positive density
martingales and fluctuation-form identities, not as an equivalence between
martingale measures and detailed balance. The Wheel's put-like payoff,
Cover's relative-regret guarantee, and information-economics results further
show what remains possible: conditional growth, benchmark-relative
performance, and transient informational rents. Any claim to work ``in all
market conditions'' must therefore disclose its benchmark, financing,
admissibility constraints, horizon, and assumptions about the market
environment.

\begin{acknowledgments}
Noson S. Yanofsky drew my attention to parallels between the present result,
Arrow's Impossibility Theorem, and fixed-point formulations of
self-reference.

This research was funded in whole or in part by the Austrian Science Fund
(FWF) [Grant DOI: 10.55776/PIN5424624]. For open access purposes, the author
has applied a CC BY public copyright license to any author accepted manuscript
version arising from this submission. The author acknowledges TU Wien
Bibliothek for financial support through its Open Access Funding Programme.

OpenAI Codex (GPT-5) was used to assist with literature organization, LaTeX
editing, and checks of derivations. The author specified the scientific
direction and prompts; model output retained in the manuscript was checked
against the cited sources and explicit calculations. The author accepts full
responsibility for the final text.
\end{acknowledgments}

\section*{Data Availability}
This is a purely mathematical work, and no data or software were created or
analyzed in this study.

%%------------------------------------------------------------------------------
\bibliography{2026-trade-additions,svozil}
\end{document}
